\chapter{Om at skrive matematik}

I \LaTeX\ skelner man mellem LR-mode (til brødtekst) og mathmode.

\section{Inline vs. display-matematik}

Matematik skrevet midt i en sætning kalder man \emph{inline}-matematik og
skrives mellem dollartegn: \texttt{\$ ... \$}. For eksempel giver
\texttt{\$f(n) = O(n \^{} 2)\$} resultatet $f(n) = O(n^2)$.
 
Hvis matematikken skal stå for sig selv på egen linje, bruges
\emph{display}-matematik med \texttt{\textbackslash[ ... \textbackslash]}
eller \texttt{equation}-omgivelsen (sidstnævnte nummererer udtrykket):
\begin{equation}
  T(n) = 2T(n/2) + n
\end{equation}
 
\section{Eksponenter, indeks og brøker}
Eksponenter skrives med \texttt{\^{}} og indeks med \texttt{\_}.
Bruges der mere end et tegn, skal det sættes i tuborgklammer \texttt{\{\}}:
 
\[
  x^2, \quad a_i, \quad 2^{n+1}, \quad x_{i,j}
\]
 
Brøker skrives med \texttt{\textbackslash frac\{tæller\}\{nævner\}}:
\[
  \frac{n!}{k!(n-k)!}
\]
 
\section{Symboler og operatorer}
\texttt{amsmath}-pakken giver adgang til mange nyttige symboler, f.eks.:
\[
  \sum_{i=1}^{n} i = \frac{n(n+1)}{2}, \qquad
  \prod_{i=1}^{n} i, \qquad
  \int_0^1 x^2 \, dx
\]
 
Mængdesymboler og logik, som ofte bruges i datalogi:
\[
  \forall x \in \mathbb{N},\ \exists y \in \mathbb{Z} : x \le y,
  \qquad A \cup B, \qquad A \cap B, \qquad \neg p \lor q
\]
 
\section{Flerlinjede udregninger}
Brug \texttt{align}-omgivelsen til at vise flere linjers udregning,
linet op efter \texttt{\&}:
\begin{align}
  (a+b)^2 &= a^2 + 2ab + b^2 \\
  &= a^2 + b^2 + 2ab
\end{align}
 
\section{Matricer}
Matricer skrives med \texttt{pmatrix} (parenteser) eller
\texttt{bmatrix} (kantede parenteser):

\[
  A = \begin{pmatrix} 1 & 2 \\ 3 & 4 \end{pmatrix}
\]
 
\section{Gode råd}
\begin{itemize}
  \item Indlæs altid pakken \texttt{amsmath} -- den udvider de
    grundlæggende matematik-kommandoer betydeligt.
  \item Brug \texttt{\textbackslash text\{...\}} til almindelig tekst
    inde i en matematik-omgivelse, f.eks. $n \text{ er lige}$.
  \item Lad være med at blande \texttt{\$...\$} og \texttt{equation}
    -- vælg den, der passer til, om udtrykket skal stå i eller uden for
    løbende tekst.
  \item Kompiler ofte, så du hurtigt opdager fejl i din matematik-syntaks.
  \end{itemize}
  
%%% Local Variables:
%%% mode: latex
%%% TeX-master: "report"
%%% End:
